ГОГОЛЬ  (падает из-за кулис на сцену и смирно лежит).

ПУШКИН (выходит, спотыкается об Гоголя и падает): Вот черт! Никак об Гоголя!

ГОГОЛЬ  (поднимаясь):  Мерзопакость какая! Отдохнуть не дадут! (Идет, спотыкается  об  Пушкина и падает).  Никак об Пушкина спотыкнулся!

ПУШКИН  (поднимаясь): Ни минуты покоя! (Идет, спотыкается об Гоголя и падает). Вот черт! Никак опять об Гоголя!

ГОГОЛЬ  (поднимаясь): Вечно во всем помеха! (Идет, спотыкается об  Пушкина и падает). Вот мерзопакость! Опять об Пушкина!

ПУШКИН  (поднимаясь):  Хулиганство! Сплошное хулиганство! (Идет,  спотыкается об Гоголя и падает). Вот черт! Опять об Гоголя!

ГОГОЛЬ  (поднимаясь): Это издевательство сплошное! (Идет, спотыкается об  Пушкина и падает). Опять об Пушкина!

ПУШКИН  (поднимаясь): Вот черт! Истинно что черт! (Идет, спотыкается об Гоголя и падает). Об Гоголя!

ГОГОЛЬ  (поднимаясь): Мерзопакость! (Идет, спотыкается об Пушкина и падает).  Об Пушкина!

ПУШКИН  (поднимаясь):   Вот   черт! (Идет, спотыкается об Гоголя и падает за кулисы). Об Гоголя!

ГОГОЛЬ  (поднимаясь): Мерзопакость! (Уходит за кулисы).

За сценой слышен голос Гоголя:

\begin{center} <<Об Пушкина!>> \end{center}
\begin{center} Занавес. \end{center}
\begin{flushright} 1934 г. \end{flushright}